\documentclass[10pt,a4paper]{amsart}
\usepackage{amsmath,amssymb,amsthm,graphicx}
% For commutative diagrams we use Xy-pic
\usepackage[all]{xy}

% We use 2cell for 2-commutative diagrams.
\xyoption{2cell}
% Theorem environments.
%
\theoremstyle{plain}
\newtheorem{theorem}[subsection]{Theorem}
\newtheorem{proposition}[subsection]{Proposition}
\newtheorem{lemma}[subsection]{Lemma}

\theoremstyle{definition}
\newtheorem{definition}[subsection]{Definition}
\newtheorem{example}[subsection]{Example}
\newtheorem{exercise}[subsection]{Exercise}
\newtheorem{situation}[subsection]{Situation}

\theoremstyle{remark}
\newtheorem{remark}[subsection]{Remark}
\newtheorem{remarks}[subsection]{Remarks}

\numberwithin{equation}{subsection}

% Macros
%
\def\lim{\mathop{\mathrm{lim}}\nolimits}
\def\colim{\mathop{\mathrm{colim}}\nolimits}
\def\Spec{\mathop{\mathrm{Spec}}}
\def\Hom{\mathop{\mathrm{Hom}}\nolimits}
\def\Ext{\mathop{\mathrm{Ext}}\nolimits}
\def\SheafHom{\mathop{\mathcal{H}\!\mathit{om}}\nolimits}
\def\SheafExt{\mathop{\mathcal{E}\!\mathit{xt}}\nolimits}
\def\Sch{\mathit{Sch}}
\def\Mor{\mathop{\mathrm{Mor}}\nolimits}
\def\Ob{\mathop{\mathrm{Ob}}\nolimits}
\def\Sh{\mathop{\mathit{Sh}}\nolimits}
\def\NL{\mathop{N\!L}\nolimits}
\def\CH{\mathop{\mathrm{CH}}\nolimits}
\def\proetale{{pro\text{-}\acute{e}tale}}
\def\etale{{\acute{e}tale}}
\def\QCoh{\mathit{QCoh}}
\def\Ker{\mathop{\mathrm{Ker}}}
\def\Im{\mathop{\mathrm{Im}}}
\def\Coker{\mathop{\mathrm{Coker}}}
\def\Coim{\mathop{\mathrm{Coim}}}

% Boxtimes
%
\DeclareMathSymbol{\boxtimes}{\mathbin}{AMSa}{"02}

%
% Macros for moduli stacks/spaces
%
\def\QCohstack{\mathcal{QC}\!\mathit{oh}}
\def\Cohstack{\mathcal{C}\!\mathit{oh}}
\def\Spacesstack{\mathcal{S}\!\mathit{paces}}
\def\Quotfunctor{\mathrm{Quot}}
\def\Hilbfunctor{\mathrm{Hilb}}
\def\Curvesstack{\mathcal{C}\!\mathit{urves}}
\def\Polarizedstack{\mathcal{P}\!\mathit{olarized}}
\def\Complexesstack{\mathcal{C}\!\mathit{omplexes}}
% \Pic is the operator that assigns to X its picard group, usage \Pic(X)
% \Picardstack_{X/B} denotes the Picard stack of X over B
% \Picardfunctor_{X/B} denotes the Picard functor of X over B
\def\Pic{\mathop{\mathrm{Pic}}\nolimits}
\def\Picardstack{\mathcal{P}\!\mathit{ic}}
\def\Picardfunctor{\mathrm{Pic}}
\def\Deformationcategory{\mathcal{D}\!\mathit{ef}}
\newenvironment{reference}{}{}
\newenvironment{slogan}{}{}
\newenvironment{history}{}{}
\usepackage{cleveref}
\crefname{originalsection}{original section}{original sections}
\Crefname{originalsection}{Original section}{Original sections}
\makeatletter
\newlabel{section-coherent-formal}{{0EHN}{source}{}{sourcepointer-1}{}}
\newlabel{section-coherent-formal@cref}{{[originalsection][0][]0EHN}{[1][0][]source}{}{}{}}
\makeatother
\crefname{originalequation}{original equation}{original equations}
\Crefname{originalequation}{Original equation}{Original equations}
\makeatletter
\newlabel{equation-completion-functor}{{0880}{source}{}{sourcepointer-2}{}}
\newlabel{equation-completion-functor@cref}{{[originalequation][0][]0880}{[1][0][]source}{}{}{}}
\makeatother
\crefname{originallemma}{original lemma}{original lemmas}
\Crefname{originallemma}{Original lemma}{Original lemmas}
\makeatletter
\newlabel{modules-lemma-internal-hom-locally-kernel-direct-sum}{{01CQ}{source}{}{sourcepointer-3}{}}
\newlabel{modules-lemma-internal-hom-locally-kernel-direct-sum@cref}{{[originallemma][0][]01CQ}{[1][0][]source}{}{}{}}
\makeatother
\crefname{originallemma}{original lemma}{original lemmas}
\Crefname{originallemma}{Original lemma}{Original lemmas}
\makeatletter
\newlabel{lemma-completion-internal-hom}{{0882}{source}{}{sourcepointer-4}{}}
\newlabel{lemma-completion-internal-hom@cref}{{[originallemma][0][]0882}{[1][0][]source}{}{}{}}
\makeatother
\crefname{originaltheorem}{original theorem}{original theorems}
\Crefname{originaltheorem}{Original theorem}{Original theorems}
\makeatletter
\newlabel{lemma-spell-out-theorem-formal-functions}{{087U}{source}{}{sourcepointer-5}{}}
\newlabel{lemma-spell-out-theorem-formal-functions@cref}{{[originaltheorem][0][]087U}{[1][0][]source}{}{}{}}
\makeatother
\crefname{originallemma}{original lemma}{original lemmas}
\Crefname{originallemma}{Original lemma}{Original lemmas}
\makeatletter
\newlabel{lemma-finite-over-rees-algebra}{{087Z}{source}{}{sourcepointer-6}{}}
\newlabel{lemma-finite-over-rees-algebra@cref}{{[originallemma][0][]087Z}{[1][0][]source}{}{}{}}
\makeatother
\crefname{originallemma}{original lemma}{original lemmas}
\Crefname{originallemma}{Original lemma}{Original lemmas}
\makeatletter
\newlabel{lemma-graded-finiteness}{{0897}{source}{}{sourcepointer-7}{}}
\newlabel{lemma-graded-finiteness@cref}{{[originallemma][0][]0897}{[1][0][]source}{}{}{}}
\makeatother
\crefname{originallemma}{original lemma}{original lemmas}
\Crefname{originallemma}{Original lemma}{Original lemmas}
\makeatletter
\newlabel{lemma-inverse-systems-surjective}{{087Y}{source}{}{sourcepointer-8}{}}
\newlabel{lemma-inverse-systems-surjective@cref}{{[originallemma][0][]087Y}{[1][0][]source}{}{}{}}
\makeatother
\crefname{originallemma}{original lemma}{original lemmas}
\Crefname{originallemma}{Original lemma}{Original lemmas}
\makeatletter
\newlabel{lemma-inverse-systems-abelian}{{087X}{source}{}{sourcepointer-9}{}}
\newlabel{lemma-inverse-systems-abelian@cref}{{[originallemma][0][]087X}{[1][0][]source}{}{}{}}
\makeatother
\crefname{originallemma}{original lemma}{original lemmas}
\Crefname{originallemma}{Original lemma}{Original lemmas}
\makeatletter
\newlabel{lemma-exact}{{0881}{source}{}{sourcepointer-10}{}}
\newlabel{lemma-exact@cref}{{[originallemma][0][]0881}{[1][0][]source}{}{}{}}
\makeatother
\crefname{originallemma}{original lemma}{original lemmas}
\Crefname{originallemma}{Original lemma}{Original lemmas}
\makeatletter
\newlabel{lemma-inverse-systems-affine}{{087W}{source}{}{sourcepointer-11}{}}
\newlabel{lemma-inverse-systems-affine@cref}{{[originallemma][0][]087W}{[1][0][]source}{}{}{}}
\makeatother
\crefname{originallemma}{original lemma}{original lemmas}
\Crefname{originallemma}{Original lemma}{Original lemmas}
\makeatletter
\newlabel{lemma-i-star-equivalence}{{087T}{source}{}{sourcepointer-12}{}}
\newlabel{lemma-i-star-equivalence@cref}{{[originallemma][0][]087T}{[1][0][]source}{}{}{}}
\makeatother
\crefname{originallemma}{original lemma}{original lemmas}
\Crefname{originallemma}{Original lemma}{Original lemmas}
\makeatletter
\newlabel{lemma-existence-easy}{{0889}{source}{}{sourcepointer-13}{}}
\newlabel{lemma-existence-easy@cref}{{[originallemma][0][]0889}{[1][0][]source}{}{}{}}
\makeatother
\crefname{originalproposition}{original proposition}{original propositions}
\Crefname{originalproposition}{Original proposition}{Original propositions}
\makeatletter
\newlabel{proposition-proper-pushforward-coherent}{{02O5}{source}{}{sourcepointer-14}{}}
\newlabel{proposition-proper-pushforward-coherent@cref}{{[originalproposition][0][]02O5}{[1][0][]source}{}{}{}}
\makeatother
\crefname{originallemma}{original lemma}{original lemmas}
\Crefname{originallemma}{Original lemma}{Original lemmas}
\makeatletter
\newlabel{lemma-ML-cohomology-powers-ideal}{{02OB}{source}{}{sourcepointer-15}{}}
\newlabel{lemma-ML-cohomology-powers-ideal@cref}{{[originallemma][0][]02OB}{[1][0][]source}{}{}{}}
\makeatother
\crefname{originaltheorem}{original theorem}{original theorems}
\Crefname{originaltheorem}{Original theorem}{Original theorems}
\makeatletter
\newlabel{theorem-formal-functions}{{02OC}{source}{}{sourcepointer-16}{}}
\newlabel{theorem-formal-functions@cref}{{[originaltheorem][0][]02OC}{[1][0][]source}{}{}{}}
\makeatother
\crefname{originallemma}{original lemma}{original lemmas}
\Crefname{originallemma}{Original lemma}{Original lemmas}
\makeatletter
\newlabel{algebra-lemma-completion-flat}{{00MB}{source}{}{sourcepointer-17}{}}
\newlabel{algebra-lemma-completion-flat@cref}{{[originallemma][0][]00MB}{[1][0][]source}{}{}{}}
\makeatother
\crefname{originallemma}{original lemma}{original lemmas}
\Crefname{originallemma}{Original lemma}{Original lemmas}
\makeatletter
\newlabel{lemma-flat-base-change-cohomology}{{02KH}{source}{}{sourcepointer-18}{}}
\newlabel{lemma-flat-base-change-cohomology@cref}{{[originallemma][0][]02KH}{[1][0][]source}{}{}{}}
\makeatother
\crefname{originallemma}{original lemma}{original lemmas}
\Crefname{originallemma}{Original lemma}{Original lemmas}
\makeatletter
\newlabel{more-algebra-lemma-neighbourhood-equivalence}{{05EC}{source}{}{sourcepointer-19}{}}
\newlabel{more-algebra-lemma-neighbourhood-equivalence@cref}{{[originallemma][0][]05EC}{[1][0][]source}{}{}{}}
\makeatother
\crefname{originallemma}{original lemma}{original lemmas}
\Crefname{originallemma}{Original lemma}{Original lemmas}
\makeatletter
\newlabel{lemma-power-ideal-kills-sheaf}{{01Y9}{source}{}{sourcepointer-20}{}}
\newlabel{lemma-power-ideal-kills-sheaf@cref}{{[originallemma][0][]01Y9}{[1][0][]source}{}{}{}}
\makeatother
\crefname{originallemma}{original lemma}{original lemmas}
\Crefname{originallemma}{Original lemma}{Original lemmas}
\makeatletter
\newlabel{lemma-acc-coherent}{{01Y8}{source}{}{sourcepointer-21}{}}
\newlabel{lemma-acc-coherent@cref}{{[originallemma][0][]01Y8}{[1][0][]source}{}{}{}}
\makeatother
\crefname{originallemma}{original lemma}{original lemmas}
\Crefname{originallemma}{Original lemma}{Original lemmas}
\makeatletter
\newlabel{lemma-chow-Noetherian}{{0200}{source}{}{sourcepointer-22}{}}
\newlabel{lemma-chow-Noetherian@cref}{{[originallemma][0][]0200}{[1][0][]source}{}{}{}}
\makeatother
\crefname{originaldefinition}{original definition}{original definitions}
\Crefname{originaldefinition}{Original definition}{Original definitions}
\makeatletter
\newlabel{schemes-definition-reduced-induced-scheme}{{01J4}{source}{}{sourcepointer-23}{}}
\newlabel{schemes-definition-reduced-induced-scheme@cref}{{[originaldefinition][0][]01J4}{[1][0][]source}{}{}{}}
\makeatother
\crefname{originaldefinition}{original definition}{original definitions}
\Crefname{originaldefinition}{Original definition}{Original definitions}
\makeatletter
\newlabel{morphisms-definition-scheme-theoretic-intersection-union}{{0C4H}{source}{}{sourcepointer-24}{}}
\newlabel{morphisms-definition-scheme-theoretic-intersection-union@cref}{{[originaldefinition][0][]0C4H}{[1][0][]source}{}{}{}}
\makeatother
\crefname{originallemma}{original lemma}{original lemmas}
\Crefname{originallemma}{Original lemma}{Original lemmas}
\makeatletter
\newlabel{morphisms-lemma-scheme-theoretic-union}{{0C4J}{source}{}{sourcepointer-25}{}}
\newlabel{morphisms-lemma-scheme-theoretic-union@cref}{{[originallemma][0][]0C4J}{[1][0][]source}{}{}{}}
\makeatother
\crefname{originallemma}{original lemma}{original lemmas}
\Crefname{originallemma}{Original lemma}{Original lemmas}
\makeatletter
\newlabel{morphisms-lemma-closed-immersion-proper}{{01W5}{source}{}{sourcepointer-26}{}}
\newlabel{morphisms-lemma-closed-immersion-proper@cref}{{[originallemma][0][]01W5}{[1][0][]source}{}{}{}}
\makeatother
\crefname{originallemma}{original lemma}{original lemmas}
\Crefname{originallemma}{Original lemma}{Original lemmas}
\makeatletter
\newlabel{morphisms-lemma-composition-proper}{{01W3}{source}{}{sourcepointer-27}{}}
\newlabel{morphisms-lemma-composition-proper@cref}{{[originallemma][0][]01W3}{[1][0][]source}{}{}{}}
\makeatother
\crefname{originallemma}{original lemma}{original lemmas}
\Crefname{originallemma}{Original lemma}{Original lemmas}
\makeatletter
\newlabel{morphisms-lemma-image-universally-closed-separated}{{09MQ}{source}{}{sourcepointer-28}{}}
\newlabel{morphisms-lemma-image-universally-closed-separated@cref}{{[originallemma][0][]09MQ}{[1][0][]source}{}{}{}}
\makeatother
\crefname{originallemma}{original lemma}{original lemmas}
\Crefname{originallemma}{Original lemma}{Original lemmas}
\makeatletter
\newlabel{morphisms-lemma-immersion-locally-finite-type}{{01T5}{source}{}{sourcepointer-29}{}}
\newlabel{morphisms-lemma-immersion-locally-finite-type@cref}{{[originallemma][0][]01T5}{[1][0][]source}{}{}{}}
\makeatother
\crefname{originallemma}{original lemma}{original lemmas}
\Crefname{originallemma}{Original lemma}{Original lemmas}
\makeatletter
\newlabel{morphisms-lemma-composition-finite-type}{{01T3}{source}{}{sourcepointer-30}{}}
\newlabel{morphisms-lemma-composition-finite-type@cref}{{[originallemma][0][]01T3}{[1][0][]source}{}{}{}}
\makeatother
\crefname{originallemma}{original lemma}{original lemmas}
\Crefname{originallemma}{Original lemma}{Original lemmas}
\makeatletter
\newlabel{morphisms-lemma-image-proper-is-proper}{{03GN}{source}{}{sourcepointer-31}{}}
\newlabel{morphisms-lemma-image-proper-is-proper@cref}{{[originallemma][0][]03GN}{[1][0][]source}{}{}{}}
\makeatother
\crefname{originallemma}{original lemma}{original lemmas}
\Crefname{originallemma}{Original lemma}{Original lemmas}
\makeatletter
\newlabel{morphisms-lemma-base-change-finite-type}{{01T4}{source}{}{sourcepointer-32}{}}
\newlabel{morphisms-lemma-base-change-finite-type@cref}{{[originallemma][0][]01T4}{[1][0][]source}{}{}{}}
\makeatother
\crefname{originaldefinition}{original definition}{original definitions}
\Crefname{originaldefinition}{Original definition}{Original definitions}
\makeatletter
\newlabel{schemes-definition-inverse-image-closed-subscheme}{{01JV}{source}{}{sourcepointer-33}{}}
\newlabel{schemes-definition-inverse-image-closed-subscheme@cref}{{[originaldefinition][0][]01JV}{[1][0][]source}{}{}{}}
\makeatother
\crefname{originallemma}{original lemma}{original lemmas}
\Crefname{originallemma}{Original lemma}{Original lemmas}
\makeatletter
\newlabel{morphisms-lemma-base-change-proper}{{01W4}{source}{}{sourcepointer-34}{}}
\newlabel{morphisms-lemma-base-change-proper@cref}{{[originallemma][0][]01W4}{[1][0][]source}{}{}{}}
\makeatother
\crefname{originallemma}{original lemma}{original lemmas}
\Crefname{originallemma}{Original lemma}{Original lemmas}
\makeatletter
\newlabel{morphisms-lemma-scheme-theoretic-image-is-proper}{{0AH6}{source}{}{sourcepointer-35}{}}
\newlabel{morphisms-lemma-scheme-theoretic-image-is-proper@cref}{{[originallemma][0][]0AH6}{[1][0][]source}{}{}{}}
\makeatother
\crefname{originallemma}{original lemma}{original lemmas}
\Crefname{originallemma}{Original lemma}{Original lemmas}
\makeatletter
\newlabel{morphisms-lemma-closed-immersion-finite}{{035C}{source}{}{sourcepointer-36}{}}
\newlabel{morphisms-lemma-closed-immersion-finite@cref}{{[originallemma][0][]035C}{[1][0][]source}{}{}{}}
\makeatother
\crefname{originallemma}{original lemma}{original lemmas}
\Crefname{originallemma}{Original lemma}{Original lemmas}
\makeatletter
\newlabel{morphisms-lemma-finite-union-finite}{{0CYI}{source}{}{sourcepointer-37}{}}
\newlabel{morphisms-lemma-finite-union-finite@cref}{{[originallemma][0][]0CYI}{[1][0][]source}{}{}{}}
\makeatother
\crefname{originallemma}{original lemma}{original lemmas}
\Crefname{originallemma}{Original lemma}{Original lemmas}
\makeatletter
\newlabel{morphisms-lemma-finite-proper}{{01WN}{source}{}{sourcepointer-38}{}}
\newlabel{morphisms-lemma-finite-proper@cref}{{[originallemma][0][]01WN}{[1][0][]source}{}{}{}}
\makeatother
\crefname{originallemma}{original lemma}{original lemmas}
\Crefname{originallemma}{Original lemma}{Original lemmas}
\makeatletter
\newlabel{morphisms-lemma-support-finite-type}{{056J}{source}{}{sourcepointer-39}{}}
\newlabel{morphisms-lemma-support-finite-type@cref}{{[originallemma][0][]056J}{[1][0][]source}{}{}{}}
\makeatother
\crefname{originaldefinition}{original definition}{original definitions}
\Crefname{originaldefinition}{Original definition}{Original definitions}
\makeatletter
\newlabel{morphisms-definition-scheme-theoretic-support}{{05JV}{source}{}{sourcepointer-40}{}}
\newlabel{morphisms-definition-scheme-theoretic-support@cref}{{[originaldefinition][0][]05JV}{[1][0][]source}{}{}{}}
\makeatother
\crefname{originallemma}{original lemma}{original lemmas}
\Crefname{originallemma}{Original lemma}{Original lemmas}
\makeatletter
\newlabel{modules-lemma-pullback-finite-type}{{01B6}{source}{}{sourcepointer-41}{}}
\newlabel{modules-lemma-pullback-finite-type@cref}{{[originallemma][0][]01B6}{[1][0][]source}{}{}{}}
\makeatother
\crefname{originaldefinition}{original definition}{original definitions}
\Crefname{originaldefinition}{Original definition}{Original definitions}
\makeatletter
\newlabel{homology-definition-serre-subcategory}{{02MO}{source}{}{sourcepointer-42}{}}
\newlabel{homology-definition-serre-subcategory@cref}{{[originaldefinition][0][]02MO}{[1][0][]source}{}{}{}}
\makeatother
\crefname{originalsection}{original section}{original sections}
\Crefname{originalsection}{Original section}{Original sections}
\makeatletter
\newlabel{schemes-section-quasi-coherent}{{01LA}{source}{}{sourcepointer-43}{}}
\newlabel{schemes-section-quasi-coherent@cref}{{[originalsection][0][]01LA}{[1][0][]source}{}{}{}}
\makeatother
\crefname{originalsection}{original section}{original sections}
\Crefname{originalsection}{Original section}{Original sections}
\makeatletter
\newlabel{modules-section-finite-type}{{01B4}{source}{}{sourcepointer-44}{}}
\newlabel{modules-section-finite-type@cref}{{[originalsection][0][]01B4}{[1][0][]source}{}{}{}}
\makeatother
\crefname{originallemma}{original lemma}{original lemmas}
\Crefname{originallemma}{Original lemma}{Original lemmas}
\makeatletter
\newlabel{morphisms-lemma-finite-type-noetherian}{{01T6}{source}{}{sourcepointer-45}{}}
\newlabel{morphisms-lemma-finite-type-noetherian@cref}{{[originallemma][0][]01T6}{[1][0][]source}{}{}{}}
\makeatother
\crefname{originalsection}{original section}{original sections}
\Crefname{originalsection}{Original section}{Original sections}
\makeatletter
\newlabel{section-coherent-sheaves}{{01XY}{source}{}{sourcepointer-46}{}}
\newlabel{section-coherent-sheaves@cref}{{[originalsection][0][]01XY}{[1][0][]source}{}{}{}}
\makeatother
\crefname{originallemma}{original lemma}{original lemmas}
\Crefname{originallemma}{Original lemma}{Original lemmas}
\makeatletter
\newlabel{lemma-finite-pushforward-coherent}{{01Y6}{source}{}{sourcepointer-47}{}}
\newlabel{lemma-finite-pushforward-coherent@cref}{{[originallemma][0][]01Y6}{[1][0][]source}{}{}{}}
\makeatother
\crefname{originallemma}{original lemma}{original lemmas}
\Crefname{originallemma}{Original lemma}{Original lemmas}
\makeatletter
\newlabel{cohomology-lemma-relative-Leray}{{01F6}{source}{}{sourcepointer-48}{}}
\newlabel{cohomology-lemma-relative-Leray@cref}{{[originallemma][0][]01F6}{[1][0][]source}{}{}{}}
\makeatother
\crefname{originallemma}{original lemma}{original lemmas}
\Crefname{originallemma}{Original lemma}{Original lemmas}
\makeatletter
\newlabel{homology-lemma-characterize-serre-subcategory}{{02MP}{source}{}{sourcepointer-49}{}}
\newlabel{homology-lemma-characterize-serre-subcategory@cref}{{[originallemma][0][]02MP}{[1][0][]source}{}{}{}}
\makeatother
\crefname{originallemma}{original lemma}{original lemmas}
\Crefname{originallemma}{Original lemma}{Original lemmas}
\makeatletter
\newlabel{morphisms-lemma-H-quasi-projective-open-H-projective}{{01WA}{source}{}{sourcepointer-50}{}}
\newlabel{morphisms-lemma-H-quasi-projective-open-H-projective@cref}{{[originallemma][0][]01WA}{[1][0][]source}{}{}{}}
\makeatother
\crefname{originallemma}{original lemma}{original lemmas}
\Crefname{originallemma}{Original lemma}{Original lemmas}
\makeatletter
\newlabel{lemma-pushforward-coherent-on-open}{{0CYJ}{source}{}{sourcepointer-51}{}}
\newlabel{lemma-pushforward-coherent-on-open@cref}{{[originallemma][0][]0CYJ}{[1][0][]source}{}{}{}}
\makeatother
\crefname{originallemma}{original lemma}{original lemmas}
\Crefname{originallemma}{Original lemma}{Original lemmas}
\makeatletter
\newlabel{algebra-lemma-radical-completion}{{05GI}{source}{}{sourcepointer-52}{}}
\newlabel{algebra-lemma-radical-completion@cref}{{[originallemma][0][]05GI}{[1][0][]source}{}{}{}}
\makeatother
\title{Grothendieck existence for coherent systems with proper support on a separated finite-type scheme}
\author{The Stacks Project Authors}
\date{stacks-project:coherent.tex:local-snapshot-20261011; GNU FDL 1.2 or later}
\begin{document}
\maketitle
\noindent\textit{Editorial note.} Public376 counts projective algebraization. This new statement allows a separated finite-type ambient X that need not be proper and requires proper support; the projective result is retained as uncounted support.\par\medskip
\noindent\textit{Editorial note.} The earlier proper case and its local lemmas are included for completeness but are not additional counted tasks.\par\medskip
\noindent\textit{Allowed standard prerequisites (editorial).}\begin{itemize}
\item Explicit prerequisite boundary. Coherent formal systems, Artin–Rees and formal functions, proper coherent pushforward, Chow lemma, Noetherian induction and support under proper maps. Projective/proper existence and the local push-pull/error and support lemmas are retained author proof, not given away as the target theorem.
\end{itemize}\medskip
\ifdefined\hypertarget\hypertarget{sourcepointer-1}{}\fi
\noindent\textit{Source reference.} Original section 0EHN: focused-seeds/AGP/sources/stacks-snapshot-20261011/tags.txt, source lines 18138--18138. Stable Stacks tag for an expressly allowed foundational or notational input; the counted construction remains the two concrete core obligations. The retained local core is not replaced by this pointer. See standard\_prerequisites.\par\medskip
\ifdefined\hypertarget\hypertarget{sourcepointer-2}{}\fi
\noindent\textit{Source reference.} Original equation 0880: focused-seeds/AGP/sources/stacks-snapshot-20261011/tags.txt, source lines 10445--10445. Stable Stacks tag for an expressly allowed foundational or notational input; the counted construction remains the two concrete core obligations. The retained local core is not replaced by this pointer. See standard\_prerequisites.\par\medskip
\ifdefined\hypertarget\hypertarget{sourcepointer-3}{}\fi
\noindent\textit{Source reference.} Original lemma 01CQ: focused-seeds/AGP/sources/stacks-snapshot-20261011/tags.txt, source lines 1771--1771. Stable Stacks tag for an expressly allowed foundational or notational input; the counted construction remains the two concrete core obligations. The retained local core is not replaced by this pointer. See standard\_prerequisites.\par\medskip
\ifdefined\hypertarget\hypertarget{sourcepointer-4}{}\fi
\noindent\textit{Source reference.} Original lemma 0882: focused-seeds/AGP/sources/stacks-snapshot-20261011/tags.txt, source lines 10447--10447. Stable Stacks tag for an expressly allowed foundational or notational input; the counted construction remains the two concrete core obligations. The retained local core is not replaced by this pointer. See standard\_prerequisites.\par\medskip
\ifdefined\hypertarget\hypertarget{sourcepointer-5}{}\fi
\noindent\textit{Source reference.} Original theorem 087U: focused-seeds/AGP/sources/stacks-snapshot-20261011/tags.txt, source lines 10439--10439. Stable Stacks tag for an expressly allowed foundational or notational input; the counted construction remains the two concrete core obligations. The retained local core is not replaced by this pointer. See standard\_prerequisites.\par\medskip
\ifdefined\hypertarget\hypertarget{sourcepointer-6}{}\fi
\noindent\textit{Source reference.} Original lemma 087Z: focused-seeds/AGP/sources/stacks-snapshot-20261011/tags.txt, source lines 10444--10444. Stable Stacks tag for an expressly allowed foundational or notational input; the counted construction remains the two concrete core obligations. The retained local core is not replaced by this pointer. See standard\_prerequisites.\par\medskip
\ifdefined\hypertarget\hypertarget{sourcepointer-7}{}\fi
\noindent\textit{Source reference.} Original lemma 0897: focused-seeds/AGP/sources/stacks-snapshot-20261011/tags.txt, source lines 10487--10487. Stable Stacks tag for an expressly allowed foundational or notational input; the counted construction remains the two concrete core obligations. The retained local core is not replaced by this pointer. See standard\_prerequisites.\par\medskip
\ifdefined\hypertarget\hypertarget{sourcepointer-8}{}\fi
\noindent\textit{Source reference.} Original lemma 087Y: focused-seeds/AGP/sources/stacks-snapshot-20261011/tags.txt, source lines 10443--10443. Stable Stacks tag for an expressly allowed foundational or notational input; the counted construction remains the two concrete core obligations. The retained local core is not replaced by this pointer. See standard\_prerequisites.\par\medskip
\ifdefined\hypertarget\hypertarget{sourcepointer-9}{}\fi
\noindent\textit{Source reference.} Original lemma 087X: focused-seeds/AGP/sources/stacks-snapshot-20261011/tags.txt, source lines 10442--10442. Stable Stacks tag for an expressly allowed foundational or notational input; the counted construction remains the two concrete core obligations. The retained local core is not replaced by this pointer. See standard\_prerequisites.\par\medskip
\ifdefined\hypertarget\hypertarget{sourcepointer-10}{}\fi
\noindent\textit{Source reference.} Original lemma 0881: focused-seeds/AGP/sources/stacks-snapshot-20261011/tags.txt, source lines 10446--10446. Stable Stacks tag for an expressly allowed foundational or notational input; the counted construction remains the two concrete core obligations. The retained local core is not replaced by this pointer. See standard\_prerequisites.\par\medskip
\ifdefined\hypertarget\hypertarget{sourcepointer-11}{}\fi
\noindent\textit{Source reference.} Original lemma 087W: focused-seeds/AGP/sources/stacks-snapshot-20261011/tags.txt, source lines 10441--10441. Stable Stacks tag for an expressly allowed foundational or notational input; the counted construction remains the two concrete core obligations. The retained local core is not replaced by this pointer. See standard\_prerequisites.\par\medskip
\ifdefined\hypertarget\hypertarget{sourcepointer-12}{}\fi
\noindent\textit{Source reference.} Original lemma 087T: focused-seeds/AGP/sources/stacks-snapshot-20261011/tags.txt, source lines 10438--10438. Stable Stacks tag for an expressly allowed foundational or notational input; the counted construction remains the two concrete core obligations. The retained local core is not replaced by this pointer. See standard\_prerequisites.\par\medskip
\ifdefined\hypertarget\hypertarget{sourcepointer-13}{}\fi
\noindent\textit{Source reference.} Original lemma 0889: focused-seeds/AGP/sources/stacks-snapshot-20261011/tags.txt, source lines 10454--10454. Stable Stacks tag for an expressly allowed foundational or notational input; the counted construction remains the two concrete core obligations. The retained local core is not replaced by this pointer. See standard\_prerequisites.\par\medskip
\ifdefined\hypertarget\hypertarget{sourcepointer-14}{}\fi
\noindent\textit{Source reference.} Original proposition 02O5: focused-seeds/AGP/sources/stacks-snapshot-20261011/tags.txt, source lines 3504--3504. Stable Stacks tag for an expressly allowed foundational or notational input; the counted construction remains the two concrete core obligations. The retained local core is not replaced by this pointer. See standard\_prerequisites.\par\medskip
\ifdefined\hypertarget\hypertarget{sourcepointer-15}{}\fi
\noindent\textit{Source reference.} Original lemma 02OB: focused-seeds/AGP/sources/stacks-snapshot-20261011/tags.txt, source lines 3510--3510. Stable Stacks tag for an expressly allowed foundational or notational input; the counted construction remains the two concrete core obligations. The retained local core is not replaced by this pointer. See standard\_prerequisites.\par\medskip
\ifdefined\hypertarget\hypertarget{sourcepointer-16}{}\fi
\noindent\textit{Source reference.} Original theorem 02OC: focused-seeds/AGP/sources/stacks-snapshot-20261011/tags.txt, source lines 3511--3511. Stable Stacks tag for an expressly allowed foundational or notational input; the counted construction remains the two concrete core obligations. The retained local core is not replaced by this pointer. See standard\_prerequisites.\par\medskip
\ifdefined\hypertarget\hypertarget{sourcepointer-17}{}\fi
\noindent\textit{Source reference.} Original lemma 00MB: focused-seeds/AGP/sources/stacks-snapshot-20261011/tags.txt, source lines 817--817. Stable Stacks tag for an expressly allowed foundational or notational input; the counted construction remains the two concrete core obligations. The retained local core is not replaced by this pointer. See standard\_prerequisites.\par\medskip
\ifdefined\hypertarget\hypertarget{sourcepointer-18}{}\fi
\noindent\textit{Source reference.} Original lemma 02KH: focused-seeds/AGP/sources/stacks-snapshot-20261011/tags.txt, source lines 3372--3372. Stable Stacks tag for an expressly allowed foundational or notational input; the counted construction remains the two concrete core obligations. The retained local core is not replaced by this pointer. See standard\_prerequisites.\par\medskip
\ifdefined\hypertarget\hypertarget{sourcepointer-19}{}\fi
\noindent\textit{Source reference.} Original lemma 05EC: focused-seeds/AGP/sources/stacks-snapshot-20261011/tags.txt, source lines 6985--6985. Stable Stacks tag for an expressly allowed foundational or notational input; the counted construction remains the two concrete core obligations. The retained local core is not replaced by this pointer. See standard\_prerequisites.\par\medskip
\ifdefined\hypertarget\hypertarget{sourcepointer-20}{}\fi
\noindent\textit{Source reference.} Original lemma 01Y9: focused-seeds/AGP/sources/stacks-snapshot-20261011/tags.txt, source lines 2548--2548. Stable Stacks tag for an expressly allowed foundational or notational input; the counted construction remains the two concrete core obligations. The retained local core is not replaced by this pointer. See standard\_prerequisites.\par\medskip
\ifdefined\hypertarget\hypertarget{sourcepointer-21}{}\fi
\noindent\textit{Source reference.} Original lemma 01Y8: focused-seeds/AGP/sources/stacks-snapshot-20261011/tags.txt, source lines 2547--2547. Stable Stacks tag for an expressly allowed foundational or notational input; the counted construction remains the two concrete core obligations. The retained local core is not replaced by this pointer. See standard\_prerequisites.\par\medskip
\ifdefined\hypertarget\hypertarget{sourcepointer-22}{}\fi
\noindent\textit{Source reference.} Original lemma 0200: focused-seeds/AGP/sources/stacks-snapshot-20261011/tags.txt, source lines 2617--2617. Stable Stacks tag for an expressly allowed foundational or notational input; the counted construction remains the two concrete core obligations. The retained local core is not replaced by this pointer. See standard\_prerequisites.\par\medskip
\ifdefined\hypertarget\hypertarget{sourcepointer-23}{}\fi
\noindent\textit{Source reference.} Original definition 01J4: focused-seeds/AGP/sources/stacks-snapshot-20261011/tags.txt, source lines 2001--2001. Stable Stacks tag for an expressly allowed foundational or notational input; the counted construction remains the two concrete core obligations. The retained local core is not replaced by this pointer. See standard\_prerequisites.\par\medskip
\ifdefined\hypertarget\hypertarget{sourcepointer-24}{}\fi
\noindent\textit{Source reference.} Original definition 0C4H: focused-seeds/AGP/sources/stacks-snapshot-20261011/tags.txt, source lines 15227--15227. Stable Stacks tag for an expressly allowed foundational or notational input; the counted construction remains the two concrete core obligations. The retained local core is not replaced by this pointer. See standard\_prerequisites.\par\medskip
\ifdefined\hypertarget\hypertarget{sourcepointer-25}{}\fi
\noindent\textit{Source reference.} Original lemma 0C4J: focused-seeds/AGP/sources/stacks-snapshot-20261011/tags.txt, source lines 15229--15229. Stable Stacks tag for an expressly allowed foundational or notational input; the counted construction remains the two concrete core obligations. The retained local core is not replaced by this pointer. See standard\_prerequisites.\par\medskip
\ifdefined\hypertarget\hypertarget{sourcepointer-26}{}\fi
\noindent\textit{Source reference.} Original lemma 01W5: focused-seeds/AGP/sources/stacks-snapshot-20261011/tags.txt, source lines 2470--2470. Stable Stacks tag for an expressly allowed foundational or notational input; the counted construction remains the two concrete core obligations. The retained local core is not replaced by this pointer. See standard\_prerequisites.\par\medskip
\ifdefined\hypertarget\hypertarget{sourcepointer-27}{}\fi
\noindent\textit{Source reference.} Original lemma 01W3: focused-seeds/AGP/sources/stacks-snapshot-20261011/tags.txt, source lines 2468--2468. Stable Stacks tag for an expressly allowed foundational or notational input; the counted construction remains the two concrete core obligations. The retained local core is not replaced by this pointer. See standard\_prerequisites.\par\medskip
\ifdefined\hypertarget\hypertarget{sourcepointer-28}{}\fi
\noindent\textit{Source reference.} Original lemma 09MQ: focused-seeds/AGP/sources/stacks-snapshot-20261011/tags.txt, source lines 12188--12188. Stable Stacks tag for an expressly allowed foundational or notational input; the counted construction remains the two concrete core obligations. The retained local core is not replaced by this pointer. See standard\_prerequisites.\par\medskip
\ifdefined\hypertarget\hypertarget{sourcepointer-29}{}\fi
\noindent\textit{Source reference.} Original lemma 01T5: focused-seeds/AGP/sources/stacks-snapshot-20261011/tags.txt, source lines 2362--2362. Stable Stacks tag for an expressly allowed foundational or notational input; the counted construction remains the two concrete core obligations. The retained local core is not replaced by this pointer. See standard\_prerequisites.\par\medskip
\ifdefined\hypertarget\hypertarget{sourcepointer-30}{}\fi
\noindent\textit{Source reference.} Original lemma 01T3: focused-seeds/AGP/sources/stacks-snapshot-20261011/tags.txt, source lines 2360--2360. Stable Stacks tag for an expressly allowed foundational or notational input; the counted construction remains the two concrete core obligations. The retained local core is not replaced by this pointer. See standard\_prerequisites.\par\medskip
\ifdefined\hypertarget\hypertarget{sourcepointer-31}{}\fi
\noindent\textit{Source reference.} Original lemma 03GN: focused-seeds/AGP/sources/stacks-snapshot-20261011/tags.txt, source lines 4534--4534. Stable Stacks tag for an expressly allowed foundational or notational input; the counted construction remains the two concrete core obligations. The retained local core is not replaced by this pointer. See standard\_prerequisites.\par\medskip
\ifdefined\hypertarget\hypertarget{sourcepointer-32}{}\fi
\noindent\textit{Source reference.} Original lemma 01T4: focused-seeds/AGP/sources/stacks-snapshot-20261011/tags.txt, source lines 2361--2361. Stable Stacks tag for an expressly allowed foundational or notational input; the counted construction remains the two concrete core obligations. The retained local core is not replaced by this pointer. See standard\_prerequisites.\par\medskip
\ifdefined\hypertarget\hypertarget{sourcepointer-33}{}\fi
\noindent\textit{Source reference.} Original definition 01JV: focused-seeds/AGP/sources/stacks-snapshot-20261011/tags.txt, source lines 2028--2028. Stable Stacks tag for an expressly allowed foundational or notational input; the counted construction remains the two concrete core obligations. The retained local core is not replaced by this pointer. See standard\_prerequisites.\par\medskip
\ifdefined\hypertarget\hypertarget{sourcepointer-34}{}\fi
\noindent\textit{Source reference.} Original lemma 01W4: focused-seeds/AGP/sources/stacks-snapshot-20261011/tags.txt, source lines 2469--2469. Stable Stacks tag for an expressly allowed foundational or notational input; the counted construction remains the two concrete core obligations. The retained local core is not replaced by this pointer. See standard\_prerequisites.\par\medskip
\ifdefined\hypertarget\hypertarget{sourcepointer-35}{}\fi
\noindent\textit{Source reference.} Original lemma 0AH6: focused-seeds/AGP/sources/stacks-snapshot-20261011/tags.txt, source lines 13219--13219. Stable Stacks tag for an expressly allowed foundational or notational input; the counted construction remains the two concrete core obligations. The retained local core is not replaced by this pointer. See standard\_prerequisites.\par\medskip
\ifdefined\hypertarget\hypertarget{sourcepointer-36}{}\fi
\noindent\textit{Source reference.} Original lemma 035C: focused-seeds/AGP/sources/stacks-snapshot-20261011/tags.txt, source lines 4123--4123. Stable Stacks tag for an expressly allowed foundational or notational input; the counted construction remains the two concrete core obligations. The retained local core is not replaced by this pointer. See standard\_prerequisites.\par\medskip
\ifdefined\hypertarget\hypertarget{sourcepointer-37}{}\fi
\noindent\textit{Source reference.} Original lemma 0CYI: focused-seeds/AGP/sources/stacks-snapshot-20261011/tags.txt, source lines 16243--16243. Stable Stacks tag for an expressly allowed foundational or notational input; the counted construction remains the two concrete core obligations. The retained local core is not replaced by this pointer. See standard\_prerequisites.\par\medskip
\ifdefined\hypertarget\hypertarget{sourcepointer-38}{}\fi
\noindent\textit{Source reference.} Original lemma 01WN: focused-seeds/AGP/sources/stacks-snapshot-20261011/tags.txt, source lines 2488--2488. Stable Stacks tag for an expressly allowed foundational or notational input; the counted construction remains the two concrete core obligations. The retained local core is not replaced by this pointer. See standard\_prerequisites.\par\medskip
\ifdefined\hypertarget\hypertarget{sourcepointer-39}{}\fi
\noindent\textit{Source reference.} Original lemma 056J: focused-seeds/AGP/sources/stacks-snapshot-20261011/tags.txt, source lines 6712--6712. Stable Stacks tag for an expressly allowed foundational or notational input; the counted construction remains the two concrete core obligations. The retained local core is not replaced by this pointer. See standard\_prerequisites.\par\medskip
\ifdefined\hypertarget\hypertarget{sourcepointer-40}{}\fi
\noindent\textit{Source reference.} Original definition 05JV: focused-seeds/AGP/sources/stacks-snapshot-20261011/tags.txt, source lines 7179--7179. Stable Stacks tag for an expressly allowed foundational or notational input; the counted construction remains the two concrete core obligations. The retained local core is not replaced by this pointer. See standard\_prerequisites.\par\medskip
\ifdefined\hypertarget\hypertarget{sourcepointer-41}{}\fi
\noindent\textit{Source reference.} Original lemma 01B6: focused-seeds/AGP/sources/stacks-snapshot-20261011/tags.txt, source lines 1715--1715. Stable Stacks tag for an expressly allowed foundational or notational input; the counted construction remains the two concrete core obligations. The retained local core is not replaced by this pointer. See standard\_prerequisites.\par\medskip
\ifdefined\hypertarget\hypertarget{sourcepointer-42}{}\fi
\noindent\textit{Source reference.} Original definition 02MO: focused-seeds/AGP/sources/stacks-snapshot-20261011/tags.txt, source lines 3451--3451. Stable Stacks tag for an expressly allowed foundational or notational input; the counted construction remains the two concrete core obligations. The retained local core is not replaced by this pointer. See standard\_prerequisites.\par\medskip
\ifdefined\hypertarget\hypertarget{sourcepointer-43}{}\fi
\noindent\textit{Source reference.} Original section 01LA: focused-seeds/AGP/sources/stacks-snapshot-20261011/tags.txt, source lines 2079--2079. Stable Stacks tag for an expressly allowed foundational or notational input; the counted construction remains the two concrete core obligations. The retained local core is not replaced by this pointer. See standard\_prerequisites.\par\medskip
\ifdefined\hypertarget\hypertarget{sourcepointer-44}{}\fi
\noindent\textit{Source reference.} Original section 01B4: focused-seeds/AGP/sources/stacks-snapshot-20261011/tags.txt, source lines 1713--1713. Stable Stacks tag for an expressly allowed foundational or notational input; the counted construction remains the two concrete core obligations. The retained local core is not replaced by this pointer. See standard\_prerequisites.\par\medskip
\ifdefined\hypertarget\hypertarget{sourcepointer-45}{}\fi
\noindent\textit{Source reference.} Original lemma 01T6: focused-seeds/AGP/sources/stacks-snapshot-20261011/tags.txt, source lines 2363--2363. Stable Stacks tag for an expressly allowed foundational or notational input; the counted construction remains the two concrete core obligations. The retained local core is not replaced by this pointer. See standard\_prerequisites.\par\medskip
\ifdefined\hypertarget\hypertarget{sourcepointer-46}{}\fi
\noindent\textit{Source reference.} Original section 01XY: focused-seeds/AGP/sources/stacks-snapshot-20261011/tags.txt, source lines 2537--2537. Stable Stacks tag for an expressly allowed foundational or notational input; the counted construction remains the two concrete core obligations. The retained local core is not replaced by this pointer. See standard\_prerequisites.\par\medskip
\ifdefined\hypertarget\hypertarget{sourcepointer-47}{}\fi
\noindent\textit{Source reference.} Original lemma 01Y6: focused-seeds/AGP/sources/stacks-snapshot-20261011/tags.txt, source lines 2545--2545. Stable Stacks tag for an expressly allowed foundational or notational input; the counted construction remains the two concrete core obligations. The retained local core is not replaced by this pointer. See standard\_prerequisites.\par\medskip
\ifdefined\hypertarget\hypertarget{sourcepointer-48}{}\fi
\noindent\textit{Source reference.} Original lemma 01F6: focused-seeds/AGP/sources/stacks-snapshot-20261011/tags.txt, source lines 1859--1859. Stable Stacks tag for an expressly allowed foundational or notational input; the counted construction remains the two concrete core obligations. The retained local core is not replaced by this pointer. See standard\_prerequisites.\par\medskip
\ifdefined\hypertarget\hypertarget{sourcepointer-49}{}\fi
\noindent\textit{Source reference.} Original lemma 02MP: focused-seeds/AGP/sources/stacks-snapshot-20261011/tags.txt, source lines 3452--3452. Stable Stacks tag for an expressly allowed foundational or notational input; the counted construction remains the two concrete core obligations. The retained local core is not replaced by this pointer. See standard\_prerequisites.\par\medskip
\ifdefined\hypertarget\hypertarget{sourcepointer-50}{}\fi
\noindent\textit{Source reference.} Original lemma 01WA: focused-seeds/AGP/sources/stacks-snapshot-20261011/tags.txt, source lines 2475--2475. Stable Stacks tag for an expressly allowed foundational or notational input; the counted construction remains the two concrete core obligations. The retained local core is not replaced by this pointer. See standard\_prerequisites.\par\medskip
\ifdefined\hypertarget\hypertarget{sourcepointer-51}{}\fi
\noindent\textit{Source reference.} Original lemma 0CYJ: focused-seeds/AGP/sources/stacks-snapshot-20261011/tags.txt, source lines 16244--16244. Stable Stacks tag for an expressly allowed foundational or notational input; the counted construction remains the two concrete core obligations. The retained local core is not replaced by this pointer. See standard\_prerequisites.\par\medskip
\ifdefined\hypertarget\hypertarget{sourcepointer-52}{}\fi
\noindent\textit{Source reference.} Original lemma 05GI: focused-seeds/AGP/sources/stacks-snapshot-20261011/tags.txt, source lines 7062--7062. Stable Stacks tag for an expressly allowed foundational or notational input; the counted construction remains the two concrete core obligations. The retained local core is not replaced by this pointer. See standard\_prerequisites.\par\medskip
\begin{theorem}[Grothendieck's existence theorem]
\label{theorem-grothendieck-existence}
\begin{reference}
\cite[III Theorem 5.1.5]{EGA}
\end{reference}
Let $A$ be a Noetherian ring complete with respect to an ideal $I$.
Let $X$ be a separated, finite type scheme over $A$. Then
the functor
(\ref{equation-completion-functor-proper-over-A})
$$
\textit{Coh}_{\text{support proper over }A}(\mathcal{O}_X)
\longrightarrow
\textit{Coh}_{\text{support proper over }A}(X, \mathcal{I})
$$
is an equivalence.
\end{theorem}

\section{Grothendieck's existence theorem, I}
\label{section-existence}

\noindent
In this section we discuss Grothendieck's existence theorem for the
projective case. We will use the notion of coherent formal modules
developed in Section \ref{section-coherent-formal}. The reader who is familiar
with formal schemes is encouraged to read the statement and proof
of the theorem in \cite{EGA}.

\begin{lemma}
\label{lemma-fully-faithful}
Let $A$ be Noetherian ring complete with respect to an ideal $I$.
Let $f : X \to \Spec(A)$ be a proper morphism. Let
$\mathcal{I} = I\mathcal{O}_X$.
Then the functor (\ref{equation-completion-functor}) is fully faithful.
\end{lemma}

\begin{proof}
Let $\mathcal{F}$, $\mathcal{G}$ be coherent $\mathcal{O}_X$-modules.
Then $\mathcal{H} = \SheafHom_{\mathcal{O}_X}(\mathcal{G}, \mathcal{F})$
is a coherent $\mathcal{O}_X$-module, see
Modules, Lemma \ref{modules-lemma-internal-hom-locally-kernel-direct-sum}.
By Lemma \ref{lemma-completion-internal-hom} the map
$$
\lim_n H^0(X, \mathcal{H}/\mathcal{I}^n\mathcal{H})
\to
\Mor_{\textit{Coh}(X, \mathcal{I})}
(\mathcal{G}^\wedge, \mathcal{F}^\wedge)
$$
is bijective. Hence fully faithfulness of
(\ref{equation-completion-functor}) follows from the theorem on formal
functions (Lemma \ref{lemma-spell-out-theorem-formal-functions})
for the coherent sheaf $\mathcal{H}$.
\end{proof}

\begin{lemma}
\label{lemma-vanishing-projective}
Let $A$ be Noetherian ring and $I \subset A$ an ideal.
Let $f : X \to \Spec(A)$ be a proper morphism and let
$\mathcal{L}$ be an $f$-ample invertible sheaf. Let
$\mathcal{I} = I\mathcal{O}_X$. Let $(\mathcal{F}_n)$ be an
object of $\textit{Coh}(X, \mathcal{I})$. Then there exists an
integer $d_0$ such that
$$
H^1(X, \Ker(\mathcal{F}_{n + 1} \to \mathcal{F}_n)
\otimes \mathcal{L}^{\otimes d} )
= 0
$$
for all $n \geq 0$ and all $d \geq d_0$.
\end{lemma}

\begin{proof}
Set $B = \bigoplus I^n/I^{n + 1}$ and
$\mathcal{B} = \bigoplus \mathcal{I}^n/\mathcal{I}^{n + 1} = f^*\widetilde{B}$.
By Lemma \ref{lemma-finite-over-rees-algebra} the graded quasi-coherent
$\mathcal{B}$-module
$\mathcal{G} = \bigoplus \Ker(\mathcal{F}_{n + 1} \to \mathcal{F}_n)$
is of finite type. Hence the lemma follows from
Lemma \ref{lemma-graded-finiteness} part (2).
\end{proof}

\begin{lemma}
\label{lemma-existence-projective}
Let $A$ be Noetherian ring complete with respect to an ideal $I$.
Let $f : X \to \Spec(A)$ be a projective morphism. Let
$\mathcal{I} = I\mathcal{O}_X$.
Then the functor (\ref{equation-completion-functor}) is an equivalence.
\end{lemma}

\begin{proof}
We have already seen that (\ref{equation-completion-functor}) is
fully faithful in Lemma \ref{lemma-fully-faithful}. Thus it suffices
to show that the functor is essentially surjective.

\medskip\noindent
We first show that every object $(\mathcal{F}_n)$ of
$\textit{Coh}(X, \mathcal{I})$ is the quotient of an object
in the image of (\ref{equation-completion-functor}). 
Let $\mathcal{L}$ be an $f$-ample invertible sheaf on $X$.
Choose $d_0$ as in Lemma \ref{lemma-vanishing-projective}.
Choose a $d \geq d_0$ such that
$\mathcal{F}_1 \otimes \mathcal{L}^{\otimes d}$
is globally generated by some sections $s_{1, 1}, \ldots, s_{t, 1}$.
Since the transition maps of the system
$$
H^0(X, \mathcal{F}_{n + 1} \otimes \mathcal{L}^{\otimes d})
\longrightarrow
H^0(X, \mathcal{F}_n \otimes \mathcal{L}^{\otimes d})
$$
are surjective by the vanishing of $H^1$ we can lift
$s_{1, 1}, \ldots, s_{t, 1}$ to a compatible system of global sections
$s_{1, n}, \ldots, s_{t, n}$ of
$\mathcal{F}_n \otimes \mathcal{L}^{\otimes d}$.
These determine a compatible system of maps
$$
(s_{1, n}, \ldots, s_{t, n}) :
(\mathcal{L}^{\otimes -d})^{\oplus t} \longrightarrow \mathcal{F}_n
$$
Using Lemma \ref{lemma-inverse-systems-surjective}
we deduce that we have a surjective map
$$
\left((\mathcal{L}^{\otimes -d})^{\oplus t}\right)^\wedge
\longrightarrow
(\mathcal{F}_n)
$$
as desired.

\medskip\noindent
The result of the previous paragraph and the fact that
$\textit{Coh}(X, \mathcal{I})$ is abelian
(Lemma \ref{lemma-inverse-systems-abelian})
implies that
every object of $\textit{Coh}(X, \mathcal{I})$ is a cokernel
of a map between objects coming from $\textit{Coh}(\mathcal{O}_X)$.
As (\ref{equation-completion-functor}) is fully faithful and exact by
Lemmas \ref{lemma-fully-faithful} and \ref{lemma-exact}
we conclude.
\end{proof}








\section{Grothendieck's existence theorem, II}
\label{section-existence-proper}

\noindent
In this section we discuss Grothendieck's existence theorem in the proper case.
Before we give the statement and proof, we need to develop a bit
more theory regarding the categories $\textit{Coh}(X, \mathcal{I})$
of coherent formal modules
introduced in Section \ref{section-coherent-formal}.

\begin{remark}
\label{remark-inverse-systems-kernel-cokernel-annihilated-by}
Let $X$ be a Noetherian scheme and let
$\mathcal{I}, \mathcal{K} \subset \mathcal{O}_X$
be quasi-coherent sheaves of ideals. Let
$\alpha : (\mathcal{F}_n) \to (\mathcal{G}_n)$ be a morphism of
$\textit{Coh}(X, \mathcal{I})$.
Given an affine open $\Spec(A) = U \subset X$ with
$\mathcal{I}|_U, \mathcal{K}|_U$ corresponding to ideals $I, K \subset A$
denote $\alpha_U : M \to N$ of finite $A^\wedge$-modules which
corresponds to $\alpha|_U$ via Lemma \ref{lemma-inverse-systems-affine}.
We claim the following are equivalent
\begin{enumerate}
\item there exists an integer $t \geq 1$ such that
$\Ker(\alpha_n)$ and $\Coker(\alpha_n)$
are annihilated by $\mathcal{K}^t$ for all $n \geq 1$,
\item for any affine open $\Spec(A) = U \subset X$ as above
the modules $\Ker(\alpha_U)$ and $\Coker(\alpha_U)$
are annihilated by $K^t$ for some integer $t \geq 1$, and
\item there exists a finite affine open covering $X = \bigcup U_i$
such that the conclusion of (2) holds for $\alpha_{U_i}$.
\end{enumerate}
If these equivalent conditions hold we will say that
$\alpha$ is a
{\it map whose kernel and cokernel are annihilated by a power of
$\mathcal{K}$}.
To see the equivalence we use the following commutative algebra fact:
suppose given an exact sequence
$$
0 \to T \to M \to N \to Q \to 0
$$
of $A$-modules with $T$ and $Q$ annihilated by $K^t$ for some
ideal $K \subset A$. Then for every $f, g \in K^t$ there exists a
canonical map $"fg": N \to M$ such that $M \to N \to M$ is equal to
multiplication by $fg$. Namely, for $y \in N$ we can pick $x \in M$
mapping to $fy$ in $N$ and then we can set $"fg"(y) = gx$. Thus it is
clear that $\Ker(M/JM \to N/JN)$ and $\Coker(M/JM \to N/JN)$
are annihilated by $K^{2t}$ for any ideal $J \subset A$.

\medskip\noindent
Applying the commutative algebra fact to $\alpha_{U_i}$ and $J = I^n$
we see that (3) implies (1). Conversely,
suppose (1) holds and $M \to N$ is equal to $\alpha_U$. Then there is
a $t \geq 1$ such that
$\Ker(M/I^nM \to N/I^nN)$ and $\Coker(M/I^nM \to N/I^nN)$
are annihilated by $K^t$ for all $n$. We obtain maps
$"fg" : N/I^nN \to M/I^nM$ which in the limit induce a map $N \to M$
as $N$ and $M$ are $I$-adically complete. Since the composition with
$N \to M \to N$ is multiplication by $fg$ we conclude that $fg$
annihilates $T$ and $Q$. In other words $T$ and $Q$ are annihilated by
$K^{2t}$ as desired.
\end{remark}

\begin{lemma}
\label{lemma-existence-tricky}
Let $X$ be a Noetherian scheme. Let
$\mathcal{I}, \mathcal{K} \subset \mathcal{O}_X$
be quasi-coherent sheaves of ideals.
Let $X_e \subset X$ be the closed subscheme cut out by $\mathcal{K}^e$.
Let $\mathcal{I}_e = \mathcal{I}\mathcal{O}_{X_e}$.
Let $(\mathcal{F}_n)$ be an object of $\textit{Coh}(X, \mathcal{I})$.
Assume
\begin{enumerate}
\item the functor
$\textit{Coh}(\mathcal{O}_{X_e}) \to \textit{Coh}(X_e, \mathcal{I}_e)$
is an equivalence for all $e \geq 1$, and
\item there exists a coherent sheaf $\mathcal{H}$ on $X$ and a map
$\alpha : (\mathcal{F}_n) \to \mathcal{H}^\wedge$ whose
kernel and cokernel are annihilated by a power of $\mathcal{K}$.
\end{enumerate}
Then $(\mathcal{F}_n)$ is in the essential image of
(\ref{equation-completion-functor}).
\end{lemma}

\begin{proof}
During this proof we will use without further mention that for a closed
immersion $i : Z \to X$ the functor $i_*$ gives an equivalence between the
category of coherent modules on $Z$ and coherent modules on $X$ annihilated
by the ideal sheaf of $Z$, see Lemma \ref{lemma-i-star-equivalence}.
In particular we may identify $\textit{Coh}(\mathcal{O}_{X_e})$
with the category of coherent $\mathcal{O}_X$-modules annihilated by
$\mathcal{K}^e$ and $\textit{Coh}(X_e, \mathcal{I}_e)$ as the full subcategory
of $\textit{Coh}(X, \mathcal{I})$ of objects annihilated by $\mathcal{K}^e$.
Moreover (1) tells us these two categories are equivalent under the
completion functor (\ref{equation-completion-functor}).

\medskip\noindent
Applying this equivalence we get a coherent $\mathcal{O}_X$-module
$\mathcal{G}_e$ annihilated by $\mathcal{K}^e$ corresponding to the system
$(\mathcal{F}_n/\mathcal{K}^e\mathcal{F}_n)$ of
$\textit{Coh}(X, \mathcal{I})$. The maps
$\mathcal{F}_n/\mathcal{K}^{e + 1}\mathcal{F}_n \to
\mathcal{F}_n/\mathcal{K}^e\mathcal{F}_n$ correspond to canonical maps
$\mathcal{G}_{e + 1} \to \mathcal{G}_e$ which induce isomorphisms
$\mathcal{G}_{e + 1}/\mathcal{K}^e\mathcal{G}_{e + 1} \to \mathcal{G}_e$.
Hence $(\mathcal{G}_e)$ is an object of $\textit{Coh}(X, \mathcal{K})$.
The map $\alpha$ induces a system of maps
$$
\mathcal{F}_n/\mathcal{K}^e\mathcal{F}_n
\longrightarrow
\mathcal{H}/(\mathcal{I}^n + \mathcal{K}^e)\mathcal{H}
$$
whence maps $\mathcal{G}_e \to \mathcal{H}/\mathcal{K}^e\mathcal{H}$
(by the equivalence of categories again).
Let $t \geq 1$ be an integer, which exists by assumption (2),
such that $\mathcal{K}^t$ annihilates the kernel and cokernel of all the maps
$\mathcal{F}_n \to \mathcal{H}/\mathcal{I}^n\mathcal{H}$.
Then $\mathcal{K}^{2t}$ annihilates the kernel and cokernel of the maps
$\mathcal{F}_n/\mathcal{K}^e\mathcal{F}_n \to
\mathcal{H}/(\mathcal{I}^n + \mathcal{K}^e)\mathcal{H}$, see
Remark \ref{remark-inverse-systems-kernel-cokernel-annihilated-by}.
Whereupon we conclude that $\mathcal{K}^{4t}$ annihilates the kernel and
the cokernel of the maps
$$
\mathcal{G}_e
\longrightarrow
\mathcal{H}/\mathcal{K}^e\mathcal{H},
$$
see Remark \ref{remark-inverse-systems-kernel-cokernel-annihilated-by}.
We apply Lemma \ref{lemma-existence-easy} to obtain a coherent
$\mathcal{O}_X$-module $\mathcal{F}$, a map
$a : \mathcal{F} \to \mathcal{H}$ and an isomorphism
$\beta : (\mathcal{G}_e) \to (\mathcal{F}/\mathcal{K}^e\mathcal{F})$
in $\textit{Coh}(X, \mathcal{K})$. Working backwards, for a given $n$
the triple
$(\mathcal{F}/\mathcal{I}^n\mathcal{F}, a \bmod \mathcal{I}^n, \beta
\bmod \mathcal{I}^n)$ is a triple as in the lemma for the morphism
$\alpha_n \bmod \mathcal{K}^e :
(\mathcal{F}_n/\mathcal{K}^e\mathcal{F}_n) \to
(\mathcal{H}/(\mathcal{I}^n + \mathcal{K}^e)\mathcal{H})$
of $\textit{Coh}(X, \mathcal{K})$. Thus the uniqueness in
Lemma \ref{lemma-existence-easy}
gives a canonical isomorphism
$\mathcal{F}/\mathcal{I}^n\mathcal{F} \to \mathcal{F}_n$
compatible with all the morphisms in sight. This finishes the proof
of the lemma.
\end{proof}

\begin{lemma}
\label{lemma-inverse-systems-push-pull}
Let $Y$ be a Noetherian scheme. Let
$\mathcal{J}, \mathcal{K} \subset \mathcal{O}_Y$
be quasi-coherent sheaves of ideals.
Let $f : X \to Y$ be a proper morphism which is an isomorphism
over $V = Y \setminus V(\mathcal{K})$.
Set $\mathcal{I} = f^{-1}\mathcal{J} \mathcal{O}_X$.
Let $(\mathcal{G}_n)$ be an object of $\textit{Coh}(Y, \mathcal{J})$,
let $\mathcal{F}$ be a coherent $\mathcal{O}_X$-module, and let
$\beta : (f^*\mathcal{G}_n)  \to \mathcal{F}^\wedge$ be an isomorphism in
$\textit{Coh}(X, \mathcal{I})$. Then there exists a map
$$
\alpha :
(\mathcal{G}_n)
\longrightarrow
(f_*\mathcal{F})^\wedge
$$
in $\textit{Coh}(Y, \mathcal{J})$ whose kernel and cokernel
are annihilated by a power of $\mathcal{K}$.
\end{lemma}

\begin{proof}
Since $f$ is a proper morphism we see that $f_*\mathcal{F}$
is a coherent $\mathcal{O}_Y$-module
(Proposition \ref{proposition-proper-pushforward-coherent}).
Thus the statement of the lemma makes sense.
Consider the compositions
$$
\gamma_n : \mathcal{G}_n \to
f_*f^*\mathcal{G}_n \to
f_*(\mathcal{F}/\mathcal{I}^n\mathcal{F}).
$$
Here the first map is the adjunction map and the second is $f_*\beta_n$.
We claim that there exists a unique $\alpha$ as in the lemma
such that the compositions
$$
\mathcal{G}_n \xrightarrow{\alpha_n}
f_*\mathcal{F}/\mathcal{J}^nf_*\mathcal{F} \to
f_*(\mathcal{F}/\mathcal{I}^n\mathcal{F})
$$
equal $\gamma_n$ for all $n$. Because of the uniqueness we may assume
that $Y = \Spec(B)$ is affine. Let $J \subset B$ corresponds to the
ideal $\mathcal{J}$. Set
$$
M_n = H^0(X, \mathcal{F}/\mathcal{I}^n\mathcal{F})
\quad\text{and}\quad
M = H^0(X, \mathcal{F})
$$
By Lemma \ref{lemma-ML-cohomology-powers-ideal} and
Theorem \ref{theorem-formal-functions}
the inverse limit of the modules
$M_n$ equals the completion $M^\wedge = \lim M/J^nM$.
Set $N_n = H^0(Y, \mathcal{G}_n)$ and $N = \lim N_n$.
Via the equivalence of categories of
Lemma \ref{lemma-inverse-systems-affine}
the finite $B^\wedge$ modules $N$ and $M^\wedge$ correspond
to $(\mathcal{G}_n)$ and $f_*\mathcal{F}^\wedge$.
It follows from this that $\alpha$ has to be the morphism of
$\textit{Coh}(Y, \mathcal{J})$ corresponding to the homomorphism
$$
\lim \gamma_n : N = \lim_n N_n \longrightarrow \lim M_n = M^\wedge
$$
of finite $B^\wedge$-modules.

\medskip\noindent
We still have to show that the kernel and cokernel of $\alpha$ are
annihilated by a power of $\mathcal{K}$. Set $Y' = \Spec(B^\wedge)$
and $X' = Y' \times_Y X$. Let $\mathcal{K}'$, $\mathcal{J}'$, $\mathcal{G}'_n$
and $\mathcal{I}'$, $\mathcal{F}'$ be the pullback of
$\mathcal{K}$, $\mathcal{J}$, $\mathcal{G}_n$ and
$\mathcal{I}$, $\mathcal{F}$, to $Y'$ and $X'$.
The projection morphism $f' : X' \to Y'$ is the base change of
$f$ by $Y' \to Y$. Note that $Y' \to Y$ is a flat morphism of schemes
as $B \to B^\wedge$ is flat by
Algebra, Lemma \ref{algebra-lemma-completion-flat}.
Hence $f'_*\mathcal{F}'$, resp.\ $f'_*(f')^*\mathcal{G}_n'$
is the pullback of $f_*\mathcal{F}$, resp.\ $f_*f^*\mathcal{G}_n$
to $Y'$ by Lemma \ref{lemma-flat-base-change-cohomology}.
The uniqueness of our construction shows the pullback of $\alpha$ to $Y'$
is the corresponding map $\alpha'$ constructed for the situation on $Y'$.
Moreover, to check that the kernel and cokernel of $\alpha$ are
annihilated by $\mathcal{K}^t$ it suffices to check that the
kernel and cokernel of $\alpha'$ are annihilated by
$(\mathcal{K}')^t$. Namely, to see this we need to check this for
kernels and cokernels of the maps $\alpha_n$ and $\alpha'_n$
(see Remark \ref{remark-inverse-systems-kernel-cokernel-annihilated-by})
and the ring map $B \to B^\wedge$ induces
an equivalence of categories between modules annihilated by
$J^n$ and $(J')^n$, see
More on Algebra, Lemma \ref{more-algebra-lemma-neighbourhood-equivalence}.
Thus we may assume $B$ is complete with respect to $J$.

\medskip\noindent
Assume $Y = \Spec(B)$ is affine, $\mathcal{J}$ corresponds to the ideal
$J \subset B$, and $B$ is complete with respect to $J$.
In this case $(\mathcal{G}_n)$ is in the essential image of the functor
$\textit{Coh}(\mathcal{O}_Y) \to \textit{Coh}(Y, \mathcal{J})$.
Say $\mathcal{G}$ is a coherent $\mathcal{O}_Y$-module such that
$(\mathcal{G}_n) = \mathcal{G}^\wedge$. Note that
$f^*(\mathcal{G}^\wedge) = (f^*\mathcal{G})^\wedge$. Hence
Lemma \ref{lemma-fully-faithful}
tells us that $\beta$ comes from an isomorphism
$b : f^*\mathcal{G} \to \mathcal{F}$
and $\alpha$ is the completion functor applied to
$$
\mathcal{G} \to f_*f^*\mathcal{G} \cong f_*\mathcal{F}
$$
Hence we are trying to verify that the kernel and cokernel of the
adjunction map $c : \mathcal{G} \to f_*f^*\mathcal{G}$ are annihilated by
a power of $\mathcal{K}$. However, since the restriction
$f|_{f^{-1}(V)} : f^{-1}(V) \to V$ is an isomorphism
we see that $c|_V$ is an isomorphism. Thus the coherent sheaves
$\Ker(c)$ and $\Coker(c)$ are supported on $V(\mathcal{K})$
hence are annihilated by a power of $\mathcal{K}$
(Lemma \ref{lemma-power-ideal-kills-sheaf}) as desired.
\end{proof}

\noindent
The following proposition is the form of Grothendieck's existence
theorem which is most often used in practice.

\begin{proposition}
\label{proposition-existence-proper}
Let $A$ be a Noetherian ring complete with respect to an ideal $I$.
Let $f : X \to \Spec(A)$ be a proper morphism of schemes.
Set $\mathcal{I} = I\mathcal{O}_X$.
Then the functor (\ref{equation-completion-functor}) is an equivalence.
\end{proposition}

\begin{proof}
We have already seen that (\ref{equation-completion-functor}) is
fully faithful in Lemma \ref{lemma-fully-faithful}. Thus it suffices
to show that the functor is essentially surjective.

\medskip\noindent
Consider the collection $\Xi$ of quasi-coherent sheaves of ideals
$\mathcal{K} \subset \mathcal{O}_X$ such that every object
$(\mathcal{F}_n)$ annihilated by $\mathcal{K}$ is in the essential image.
We want to show $(0)$ is in $\Xi$. If not, then since $X$ is Noetherian
there exists a maximal quasi-coherent sheaf of ideals $\mathcal{K}$
not in $\Xi$, see
Lemma \ref{lemma-acc-coherent}.
After replacing $X$ by the closed subscheme of $X$
corresponding to $\mathcal{K}$ we may assume that every nonzero
$\mathcal{K}$ is in $\Xi$. (This uses the correspondence by
coherent modules annihilated by $\mathcal{K}$ and coherent modules
on the closed subscheme corresponding to $\mathcal{K}$, see
Lemma \ref{lemma-i-star-equivalence}.)
Let $(\mathcal{F}_n)$ be an object of
$\textit{Coh}(X, \mathcal{I})$.
We will show that this object is in the essential image of the
functor (\ref{equation-completion-functor}), thereby completion the
proof of the proposition.

\medskip\noindent
Apply Chow's lemma (Lemma \ref{lemma-chow-Noetherian}) to find a
proper surjective morphism $f : X' \to X$ which is an isomorphism
over a dense open $U \subset X$ such that $X'$ is projective over $A$.
Let $\mathcal{K}$ be the quasi-coherent sheaf of ideals cutting
out the reduced complement $X \setminus U$. By the projective
case of Grothendieck's existence theorem
(Lemma \ref{lemma-existence-projective})
there exists a coherent module $\mathcal{F}'$ on $X'$ such
that $(\mathcal{F}')^\wedge \cong (f^*\mathcal{F}_n)$. By
Proposition \ref{proposition-proper-pushforward-coherent}
the $\mathcal{O}_X$-module $\mathcal{H} = f_*\mathcal{F}'$ is coherent
and by Lemma \ref{lemma-inverse-systems-push-pull}
there exists a morphism $(\mathcal{F}_n) \to \mathcal{H}^\wedge$
of $\textit{Coh}(X, \mathcal{I})$ whose kernel and cokernel are
annihilated by a power of $\mathcal{K}$. The powers $\mathcal{K}^e$
are all in $\Xi$ so that (\ref{equation-completion-functor})
is an equivalence for the closed subschemes $X_e = V(\mathcal{K}^e)$.
We conclude by Lemma \ref{lemma-existence-tricky}.
\end{proof}





\section{Being proper over a base}
\label{section-proper-over-base}

\noindent
This is just a short section to point out some useful features
of closed subsets proper over a base and finite type, quasi-coherent
modules with support proper over a base.

\begin{lemma}
\label{lemma-closed-proper-over-base}
Let $f : X \to S$ be a morphism of schemes which is locally of finite type.
Let $Z \subset X$ be a closed subset. The following are equivalent
\begin{enumerate}
\item the morphism $Z \to S$ is proper if $Z$ is endowed with the reduced
induced closed subscheme structure
(Schemes, Definition \ref{schemes-definition-reduced-induced-scheme}),
\item for some closed subscheme structure on $Z$ the morphism $Z \to S$
is proper,
\item for any closed subscheme structure on $Z$ the morphism
$Z \to S$ is proper.
\end{enumerate}
\end{lemma}

\begin{proof}
The implications (3) $\Rightarrow$ (1) and (1) $\Rightarrow$ (2)
are immediate. Thus it suffices to prove that (2) implies (3).
We urge the reader to find their own proof of this fact.
Let $Z'$ and $Z''$ be closed subscheme structures on $Z$
such that $Z' \to S$ is proper. We have to show that $Z'' \to S$ is proper.
Let $Z''' = Z' \cup Z''$ be the scheme theoretic union, see
Morphisms, Definition
\ref{morphisms-definition-scheme-theoretic-intersection-union}.
Then $Z'''$ is another closed subscheme structure on $Z$.
This follows for example from the description of scheme theoretic unions in
Morphisms, Lemma \ref{morphisms-lemma-scheme-theoretic-union}.
Since $Z'' \to Z'''$ is a closed immersion it suffices to prove
that $Z''' \to S$ is proper (see
Morphisms, Lemmas \ref{morphisms-lemma-closed-immersion-proper} and
\ref{morphisms-lemma-composition-proper}).
The morphism $Z' \to Z'''$ is a bijective closed immersion
and in particular surjective and universally closed.
Then the fact that $Z' \to S$ is separated implies that
$Z''' \to S$ is separated, see
Morphisms, Lemma \ref{morphisms-lemma-image-universally-closed-separated}.
Moreover $Z''' \to S$ is locally of finite type
as $X \to S$ is locally of finite type
(Morphisms, Lemmas \ref{morphisms-lemma-immersion-locally-finite-type} and
\ref{morphisms-lemma-composition-finite-type}).
Since $Z' \to S$ is quasi-compact and $Z' \to Z'''$ is a homeomorphism
we see that $Z''' \to S$ is quasi-compact.
Finally, since $Z' \to S$ is universally closed, we see that
the same thing is true for $Z''' \to S$ by
Morphisms, Lemma \ref{morphisms-lemma-image-proper-is-proper}.
This finishes the proof.
\end{proof}

\begin{definition}
\label{definition-proper-over-base}
Let $f : X \to S$ be a morphism of schemes which is locally of finite type.
Let $Z \subset X$ be a closed subset.
We say {\it $Z$ is proper over $S$}
if the equivalent conditions of Lemma \ref{lemma-closed-proper-over-base}
are satisfied.
\end{definition}

\noindent
The lemma used in the definition above is false if the morphism
$f : X \to S$ is not locally of finite type. Therefore we urge
the reader not to use this terminology if $f$ is not locally of
finite type.

\begin{lemma}
\label{lemma-closed-closed-proper-over-base}
Let $f : X \to S$ be a morphism of schemes which is locally of finite type.
Let $Y \subset Z \subset X$ be closed subsets.
If $Z$ is proper over $S$, then the same is true for $Y$.
\end{lemma}

\begin{proof}
Omitted.
\end{proof}

\begin{lemma}
\label{lemma-base-change-closed-proper-over-base}
Consider a cartesian diagram of schemes
$$
\xymatrix{
X' \ar[d]_{f'} \ar[r]_{g'} & X \ar[d]^f \\
S' \ar[r]^g & S
}
$$
with $f$ locally of finite type.
If $Z$ is a closed subset of $X$ proper over $S$, then
$(g')^{-1}(Z)$ is a closed subset of $X'$ proper over $S'$.
\end{lemma}

\begin{proof}
Observe that the statement makes sense as $f'$ is locally of
finite type by Morphisms, Lemma \ref{morphisms-lemma-base-change-finite-type}.
Endow $Z$ with the reduced induced closed subscheme structure.
Denote $Z' = (g')^{-1}(Z)$ the scheme theoretic inverse image
(Schemes, Definition \ref{schemes-definition-inverse-image-closed-subscheme}).
Then $Z' = X' \times_X Z = (S' \times_S X) \times_X Z = S' \times_S Z$
is proper over $S'$ as a base change of $Z$ over $S$
(Morphisms, Lemma \ref{morphisms-lemma-base-change-proper}).
\end{proof}

\begin{lemma}
\label{lemma-functoriality-closed-proper-over-base}
Let $S$ be a scheme. Let $f : X \to Y$ be a morphism of schemes which
are locally of finite type over $S$.
\begin{enumerate}
\item If $Y$ is separated over $S$ and $Z \subset X$ is a closed subset
proper over $S$, then $f(Z)$ is a closed subset of $Y$ proper over $S$.
\item If $f$ is universally closed and $Z \subset X$ is a
closed subset proper over $S$, then $f(Z)$ is a closed subset
of $Y$ proper over $S$.
\item If $f$ is proper and $Z \subset Y$ is a closed subset
proper over $S$, then $f^{-1}(Z)$ is a closed subset of $X$ proper over $S$.
\end{enumerate}
\end{lemma}

\begin{proof}
Proof of (1). Assume $Y$ is separated over $S$ and $Z \subset X$
is a closed subset proper over $S$. Endow $Z$ with the reduced induced
closed subscheme structure and apply
Morphisms, Lemma \ref{morphisms-lemma-scheme-theoretic-image-is-proper}
to $Z \to Y$ over $S$ to conclude.

\medskip\noindent
Proof of (2). Assume $f$ is universally closed and $Z \subset X$ is a
closed subset proper over $S$. Endow $Z$ and $Z' = f(Z)$ with their reduced
induced closed subscheme structures. We obtain an induced
morphism $Z \to Z'$.
Denote $Z'' = f^{-1}(Z')$ the scheme theoretic inverse image
(Schemes, Definition \ref{schemes-definition-inverse-image-closed-subscheme}).
Then $Z'' \to Z'$ is universally closed as a base change of $f$
(Morphisms, Lemma \ref{morphisms-lemma-base-change-proper}).
Hence $Z \to Z'$ is universally closed as a composition of
the closed immersion $Z \to Z''$ and $Z'' \to Z'$
(Morphisms, Lemmas
\ref{morphisms-lemma-closed-immersion-proper} and
\ref{morphisms-lemma-composition-proper}).
We conclude that $Z' \to S$ is separated by
Morphisms, Lemma \ref{morphisms-lemma-image-universally-closed-separated}.
Since $Z \to S$ is quasi-compact and $Z \to Z'$ is surjective
we see that $Z' \to S$ is quasi-compact.
Since $Z' \to S$ is the composition of $Z' \to Y$ and $Y \to S$
we see that $Z' \to S$ is locally of finite type
(Morphisms, Lemmas \ref{morphisms-lemma-immersion-locally-finite-type} and
\ref{morphisms-lemma-composition-finite-type}).
Finally, since $Z \to S$ is universally closed, we see that
the same thing is true for $Z' \to S$ by
Morphisms, Lemma \ref{morphisms-lemma-image-proper-is-proper}.
This finishes the proof.

\medskip\noindent
Proof of (3). Assume $f$ is proper and $Z \subset Y$ is a closed subset
proper over $S$. Endow $Z$ with the reduced induced closed subscheme
structure. Denote $Z' = f^{-1}(Z)$ the scheme theoretic inverse image
(Schemes, Definition \ref{schemes-definition-inverse-image-closed-subscheme}).
Then $Z' \to Z$ is proper as a base change of $f$
(Morphisms, Lemma \ref{morphisms-lemma-base-change-proper}).
Whence $Z' \to S$ is proper as the composition of $Z' \to Z$
and $Z \to S$
(Morphisms, Lemma \ref{morphisms-lemma-composition-proper}).
This finishes the proof.
\end{proof}

\begin{lemma}
\label{lemma-union-closed-proper-over-base}
Let $f : X \to S$ be a morphism of schemes which is locally of finite type.
Let $Z_i \subset X$, $i = 1, \ldots, n$ be closed subsets.
If $Z_i$, $i = 1, \ldots, n$ are proper over $S$, then the same is
true for $Z_1 \cup \ldots \cup Z_n$.
\end{lemma}

\begin{proof}
Endow $Z_i$ with their reduced induced closed subscheme structures.
The morphism
$$
Z_1 \amalg \ldots \amalg Z_n \longrightarrow X
$$
is finite by Morphisms, Lemmas
\ref{morphisms-lemma-closed-immersion-finite} and
\ref{morphisms-lemma-finite-union-finite}.
As finite morphisms are universally closed
(Morphisms, Lemma \ref{morphisms-lemma-finite-proper})
and since $Z_1 \amalg \ldots \amalg Z_n$ is proper over $S$
we conclude by
Lemma \ref{lemma-functoriality-closed-proper-over-base} part (2)
that the image $Z_1 \cup \ldots \cup Z_n$ is proper over $S$.
\end{proof}

\noindent
Let $f : X \to S$ be a morphism of schemes which is locally
of finite type. Let $\mathcal{F}$ be a finite type, quasi-coherent
$\mathcal{O}_X$-module. Then the support $\text{Supp}(\mathcal{F})$
of $\mathcal{F}$ is a closed subset of $X$, see
Morphisms, Lemma \ref{morphisms-lemma-support-finite-type}.
Hence it makes sense to say
``the support of $\mathcal{F}$ is proper over $S$''.

\begin{lemma}
\label{lemma-module-support-proper-over-base}
Let $f : X \to S$ be a morphism of schemes which is locally
of finite type. Let $\mathcal{F}$ be a finite type, quasi-coherent
$\mathcal{O}_X$-module. The following are equivalent
\begin{enumerate}
\item the support of $\mathcal{F}$ is proper over $S$,
\item the scheme theoretic support of $\mathcal{F}$
(Morphisms, Definition \ref{morphisms-definition-scheme-theoretic-support})
is proper over $S$, and
\item there exists a closed subscheme $Z \subset X$ and
a finite type, quasi-coherent $\mathcal{O}_Z$-module
$\mathcal{G}$ such that (a) $Z \to S$ is proper, and (b)
$(Z \to X)_*\mathcal{G} = \mathcal{F}$.
\end{enumerate}
\end{lemma}

\begin{proof}
The support $\text{Supp}(\mathcal{F})$ of $\mathcal{F}$ is a closed subset
of $X$, see Morphisms, Lemma \ref{morphisms-lemma-support-finite-type}.
Hence we can apply Definition \ref{definition-proper-over-base}.
Since the scheme theoretic support of $\mathcal{F}$ is a closed
subscheme whose underlying closed subset is $\text{Supp}(\mathcal{F})$
we see that (1) and (2) are equivalent by
Definition \ref{definition-proper-over-base}.
It is clear that (2) implies (3).
Conversely, if (3) is true, then
$\text{Supp}(\mathcal{F}) \subset Z$
(an inclusion of closed subsets of $X$)
and hence $\text{Supp}(\mathcal{F})$
is proper over $S$ for example by
Lemma \ref{lemma-closed-closed-proper-over-base}.
\end{proof}

\begin{lemma}
\label{lemma-base-change-module-support-proper-over-base}
Consider a cartesian diagram of schemes
$$
\xymatrix{
X' \ar[d]_{f'} \ar[r]_{g'} & X \ar[d]^f \\
S' \ar[r]^g & S
}
$$
with $f$ locally of finite type. Let $\mathcal{F}$ be a
finite type, quasi-coherent $\mathcal{O}_X$-module.
If the support of $\mathcal{F}$ is proper over $S$, then
the support of $(g')^*\mathcal{F}$ is proper over $S'$.
\end{lemma}

\begin{proof}
Observe that the statement makes sense because
$(g')*\mathcal{F}$ is of finite type by
Modules, Lemma \ref{modules-lemma-pullback-finite-type}.
We have $\text{Supp}((g')^*\mathcal{F}) = (g')^{-1}(\text{Supp}(\mathcal{F}))$
by Morphisms, Lemma \ref{morphisms-lemma-support-finite-type}.
Thus the lemma follows from
Lemma \ref{lemma-base-change-closed-proper-over-base}.
\end{proof}

\begin{lemma}
\label{lemma-cat-module-support-proper-over-base}
Let $f : X \to S$ be a morphism of schemes which is locally
of finite type. Let $\mathcal{F}$, $\mathcal{G}$
be finite type, quasi-coherent $\mathcal{O}_X$-module.
\begin{enumerate}
\item If the supports of $\mathcal{F}$, $\mathcal{G}$
are proper over $S$, then the same is true
for $\mathcal{F} \oplus \mathcal{G}$, for any extension
of $\mathcal{G}$ by $\mathcal{F}$, for $\Im(u)$ and $\Coker(u)$
given any $\mathcal{O}_X$-module map $u : \mathcal{F} \to \mathcal{G}$,
and for any quasi-coherent quotient of $\mathcal{F}$ or $\mathcal{G}$.
\item If $S$ is locally Noetherian, then the category of
coherent $\mathcal{O}_X$-modules with support proper over
$S$ is a Serre subcategory (Homology, Definition
\ref{homology-definition-serre-subcategory})
of the abelian category of
coherent $\mathcal{O}_X$-modules.
\end{enumerate}
\end{lemma}

\begin{proof}
Proof of (1). Let $Z$, $Z'$ be the support of $\mathcal{F}$
and $\mathcal{G}$. Then all the sheaves mentioned in (1)
have support contained in $Z \cup Z'$. Thus the assertion itself
is clear from Lemmas \ref{lemma-closed-closed-proper-over-base} and
\ref{lemma-union-closed-proper-over-base}
provided we check that these sheaves are finite type
and quasi-coherent. For quasi-coherence we refer the reader to
Schemes, Section \ref{schemes-section-quasi-coherent}.
For ``finite type'' we suggest the reader take a look at
Modules, Section \ref{modules-section-finite-type}.

\medskip\noindent
Proof of (2). The proof is the same as the proof of (1). Note that
the assertions make sense as $X$ is locally Noetherian by
Morphisms, Lemma \ref{morphisms-lemma-finite-type-noetherian}
and by the description of the category of coherent
modules in Section \ref{section-coherent-sheaves}.
\end{proof}

\begin{lemma}
\label{lemma-support-proper-over-base-pushforward}
Let $S$ be a locally Noetherian scheme.
Let $f : X \to S$ be a morphism of schemes which is locally of finite type.
Let $\mathcal{F}$ be a coherent $\mathcal{O}_X$-module
with support proper over $S$. Then $R^pf_*\mathcal{F}$
is a coherent $\mathcal{O}_S$-module for all $p \geq 0$.
\end{lemma}

\begin{proof}
By Lemma \ref{lemma-module-support-proper-over-base}
there exists a closed immersion $i : Z \to X$ and
a finite type, quasi-coherent $\mathcal{O}_Z$-module
$\mathcal{G}$ such that (a) $g = f \circ i : Z \to S$ is proper, and (b)
$i_*\mathcal{G} = \mathcal{F}$.
We see that $R^pg_*\mathcal{G}$ is coherent on $S$ by
Proposition \ref{proposition-proper-pushforward-coherent}.
On the other hand, $R^qi_*\mathcal{G} = 0$ for $q > 0$
(Lemma \ref{lemma-finite-pushforward-coherent}).
By Cohomology, Lemma \ref{cohomology-lemma-relative-Leray}
we get $R^pf_*\mathcal{F} = R^pg_*\mathcal{G}$ which concludes the proof.
\end{proof}

\begin{lemma}
\label{lemma-systems-with-proper-support}
Let $S$ be a Noetherian scheme. Let $f : X \to S$ be a finite type morphism.
Let $\mathcal{I} \subset \mathcal{O}_X$ be
a quasi-coherent sheaf of ideals. The following are Serre subcategories
of $\textit{Coh}(X, \mathcal{I})$
\begin{enumerate}
\item the full subcategory of $\textit{Coh}(X, \mathcal{I})$
consisting of those objects $(\mathcal{F}_n)$ such that
the support of $\mathcal{F}_1$ is proper over $S$,
\item the full subcategory of $\textit{Coh}(X, \mathcal{I})$
consisting of those objects $(\mathcal{F}_n)$ such that
there exists a closed subscheme $Z \subset X$ proper over $S$
with $\mathcal{I}_Z \mathcal{F}_n = 0$ for all $n \geq 1$.
\end{enumerate}
\end{lemma}

\begin{proof}
We will use the criterion of
Homology, Lemma \ref{homology-lemma-characterize-serre-subcategory}.
Moreover, we will use that if
$0 \to (\mathcal{G}_n) \to (\mathcal{F}_n) \to (\mathcal{H}_n) \to 0$
is a short exact sequence of $\textit{Coh}(X, \mathcal{I})$, then
(a) $\mathcal{G}_n \to \mathcal{F}_n \to \mathcal{H}_n \to 0$
is exact for all $n \geq 1$ and
(b) $\mathcal{G}_n$ is a quotient of $\Ker(\mathcal{F}_m \to \mathcal{H}_m)$
for some $m \geq n$. See proof of Lemma \ref{lemma-inverse-systems-abelian}.

\medskip\noindent
Proof of (1). Let $(\mathcal{F}_n)$ be an object of
$\textit{Coh}(X, \mathcal{I})$. Then
$\text{Supp}(\mathcal{F}_n) = \text{Supp}(\mathcal{F}_1)$ for all $n \geq 1$.
Hence by remarks (a) and (b) above we see that
for any short exact sequence
$0 \to (\mathcal{G}_n) \to (\mathcal{F}_n) \to (\mathcal{H}_n) \to 0$
of $\textit{Coh}(X, \mathcal{I})$ we have
$\text{Supp}(\mathcal{G}_1) \cup \text{Supp}(\mathcal{H}_1) =
\text{Supp}(\mathcal{F}_1)$.
This proves that the category defined in (1)
is a Serre subcategory of $\textit{Coh}(X, \mathcal{I})$.

\medskip\noindent
Proof of (2). Here we argue the same way. Let
$0 \to (\mathcal{G}_n) \to (\mathcal{F}_n) \to (\mathcal{H}_n) \to 0$
be a short exact sequence of $\textit{Coh}(X, \mathcal{I})$.
If $Z \subset X$ is a closed subscheme and $\mathcal{I}_Z$
annihilates $\mathcal{F}_n$ for all $n$, then
$\mathcal{I}_Z$ annihilates $\mathcal{G}_n$ and $\mathcal{H}_n$
for all $n$ by (a) and (b) above.
Hence if $Z \to S$ is proper, then we conclude that the category
defined in (2) is closed under taking sub and quotient objects
inside of $\textit{Coh}(X, \mathcal{I})$.
Finally, suppose that $Z \subset X$ and $Y \subset X$ are
closed subschemes proper over $S$ such that
$\mathcal{I}_Z \mathcal{G}_n = 0$ and
$\mathcal{I}_Y \mathcal{H}_n = 0$ for all $n \geq 1$.
Then it follows from (a) above that
$\mathcal{I}_{Z \cup Y} = \mathcal{I}_Z \cdot \mathcal{I}_Y$
annihilates $\mathcal{F}_n$ for all $n$.
By Lemma \ref{lemma-union-closed-proper-over-base}
(and via Definition \ref{definition-proper-over-base} which
tells us we may choose an arbitrary scheme structure used on the union)
we see that $Z \cup Y \to S$ is proper and the proof is complete.
\end{proof}








\section{Grothendieck's existence theorem, III}
\label{section-existence-proper-support}

\noindent
To state the general version of Grothendieck's existence theorem
we introduce a bit more notation. Let $A$ be a Noetherian ring
complete with respect to an ideal $I$. Let $f : X \to \Spec(A)$
be a separated finite type morphism of schemes. Set
$\mathcal{I} = I\mathcal{O}_X$. In this situation we let
$$
\textit{Coh}_{\text{support proper over } A}(\mathcal{O}_X)
$$
be the full subcategory of $\textit{Coh}(\mathcal{O}_X)$
consisting of those coherent $\mathcal{O}_X$-modules whose
support is proper over $\Spec(A)$. This is a Serre subcategory of
$\textit{Coh}(\mathcal{O}_X)$, see
Lemma \ref{lemma-cat-module-support-proper-over-base}.
Similarly, we let
$$
\textit{Coh}_{\text{support proper over } A}(X, \mathcal{I})
$$
be the full subcategory of $\textit{Coh}(X, \mathcal{I})$
consisting of those objects $(\mathcal{F}_n)$ such that
the support of $\mathcal{F}_1$ is proper over $\Spec(A)$.
This is a Serre subcategory of $\textit{Coh}(X, \mathcal{I})$
by Lemma \ref{lemma-systems-with-proper-support} part (1).
Since the support of a quotient module is contained in the support
of the module, it follows that (\ref{equation-completion-functor})
induces a functor
\begin{equation}
\label{equation-completion-functor-proper-over-A}
\textit{Coh}_{\text{support proper over }A}(\mathcal{O}_X)
\longrightarrow
\textit{Coh}_{\text{support proper over }A}(X, \mathcal{I})
\end{equation}
We are now ready to state the main theorem of this section.

\begin{proof}
We will use the equivalence of categories of
Lemma \ref{lemma-i-star-equivalence}
without further mention.
For a closed subscheme $Z \subset X$ proper over $A$
in this proof we will say a coherent module on $X$ is
``supported on $Z$'' if it is annihilated by the ideal
sheaf of $Z$ or equivalently if it is the pushforward
of a coherent module on $Z$.
By Proposition \ref{proposition-existence-proper} we know
that the result is true for
the functor between coherent modules and systems of coherent
modules supported on $Z$. Hence it suffices to show that
every object of
$\textit{Coh}_{\text{support proper over }A}(\mathcal{O}_X)$
and every object of
$\textit{Coh}_{\text{support proper over }A}(X, \mathcal{I})$ is
supported on a closed subscheme $Z \subset X$ proper over $A$.
This holds by definition for objects of
$\textit{Coh}_{\text{support proper over }A}(\mathcal{O}_X)$.
We will prove this statement for objects of
$\textit{Coh}_{\text{support proper over }A}(X, \mathcal{I})$
using the method of proof of Proposition \ref{proposition-existence-proper}.
We urge the reader to read that proof first.

\medskip\noindent
Consider the collection $\Xi$ of quasi-coherent sheaves of ideals
$\mathcal{K} \subset \mathcal{O}_X$ such that the statement holds
for every object $(\mathcal{F}_n)$ of
$\textit{Coh}_{\text{support proper over }A}(X, \mathcal{I})$
annihilated by $\mathcal{K}$. We want to show $(0)$ is in $\Xi$.
If not, then since $X$ is Noetherian there exists a maximal
quasi-coherent sheaf of ideals $\mathcal{K}$ not in $\Xi$, see
Lemma \ref{lemma-acc-coherent}.
After replacing $X$ by the closed subscheme of $X$
corresponding to $\mathcal{K}$ we may assume that every nonzero
$\mathcal{K}$ is in $\Xi$. Let $(\mathcal{F}_n)$ be an object of
$\textit{Coh}_{\text{support proper over }A}(X, \mathcal{I})$.
We will show that this object is supported on a closed subscheme
$Z \subset X$ proper over $A$, thereby completing the
proof of the theorem.

\medskip\noindent
Apply Chow's lemma (Lemma \ref{lemma-chow-Noetherian}) to find a
proper surjective morphism $f : Y \to X$ which is an isomorphism
over a dense open $U \subset X$ such that $Y$ is H-quasi-projective
over $A$. Choose an open immersion $j : Y \to Y'$ with
$Y'$ projective over $A$, see
Morphisms, Lemma \ref{morphisms-lemma-H-quasi-projective-open-H-projective}.
Observe that
$$
\text{Supp}(f^*\mathcal{F}_n) = f^{-1}\text{Supp}(\mathcal{F}_n) =
f^{-1}\text{Supp}(\mathcal{F}_1)
$$
The first equality by
Morphisms, Lemma \ref{morphisms-lemma-support-finite-type}.
By assumption and
Lemma \ref{lemma-functoriality-closed-proper-over-base} part (3)
we see that $f^{-1}\text{Supp}(\mathcal{F}_1)$ is proper over $A$.
Hence the image of $f^{-1}\text{Supp}(\mathcal{F}_1)$
under $j$ is closed in $Y'$ by
Lemma \ref{lemma-functoriality-closed-proper-over-base} part (1).
Thus $\mathcal{F}'_n = j_*f^*\mathcal{F}_n$ is coherent on
$Y'$ by Lemma \ref{lemma-pushforward-coherent-on-open}.
It follows that $(\mathcal{F}_n')$
is an object of $\textit{Coh}(Y', I\mathcal{O}_{Y'})$.
By the projective case of Grothendieck's existence theorem
(Lemma \ref{lemma-existence-projective})
there exists a coherent $\mathcal{O}_{Y'}$-module
$\mathcal{F}'$ and an isomorphism
$(\mathcal{F}')^\wedge \cong (\mathcal{F}'_n)$ in
$\textit{Coh}(Y', I\mathcal{O}_{Y'})$.
Since $\mathcal{F}'/I\mathcal{F}' = \mathcal{F}'_1$ we see that
$$
\text{Supp}(\mathcal{F}') \cap V(I\mathcal{O}_{Y'}) =
\text{Supp}(\mathcal{F}'_1) = j(f^{-1}\text{Supp}(\mathcal{F}_1))
$$
The structure morphism $p' : Y' \to \Spec(A)$ is proper, hence
$p'(\text{Supp}(\mathcal{F}') \setminus j(Y))$
is closed in $\Spec(A)$. A nonempty closed subset of $\Spec(A)$
contains a point of $V(I)$ as $I$ is contained in the Jacobson radical
of $A$ by Algebra, Lemma \ref{algebra-lemma-radical-completion}.
The displayed equation shows that
$\text{Supp}(\mathcal{F}') \cap (p')^{-1}V(I) \subset j(Y)$
hence we conclude that $\text{Supp}(\mathcal{F}') \subset j(Y)$.
Thus $\mathcal{F}'|_Y = j^*\mathcal{F}'$
is supported on a closed subscheme $Z'$ of $Y$ proper over $A$
and $(\mathcal{F}'|_Y)^\wedge = (f^*\mathcal{F}_n)$.

\medskip\noindent
Let $\mathcal{K}$ be the quasi-coherent sheaf of ideals cutting
out the reduced complement $X \setminus U$. By
Proposition \ref{proposition-proper-pushforward-coherent}
the $\mathcal{O}_X$-module $\mathcal{H} = f_*(\mathcal{F}'|_Y)$ is coherent
and by Lemma \ref{lemma-inverse-systems-push-pull}
there exists a morphism $\alpha : (\mathcal{F}_n) \to \mathcal{H}^\wedge$
of $\textit{Coh}(X, \mathcal{I})$ whose kernel and cokernel are
annihilated by a power $\mathcal{K}^t$ of $\mathcal{K}$.
We obtain an exact sequence
$$
0 \to \Ker(\alpha) \to (\mathcal{F}_n) \to
\mathcal{H}^\wedge \to \Coker(\alpha) \to 0
$$
in $\textit{Coh}(X, \mathcal{I})$. If $Z_0 \subset X$ is the scheme theoretic
support of $\mathcal{H}$, then it is clear that $Z_0 \subset f(Z')$
set-theoretically. Hence $Z_0$ is proper over $A$ by
Lemma \ref{lemma-closed-closed-proper-over-base} and
Lemma \ref{lemma-functoriality-closed-proper-over-base} part (2).
Hence $\mathcal{H}^\wedge$ is in the subcategory defined in
Lemma \ref{lemma-systems-with-proper-support} part (2)
and a fortiori in
$\textit{Coh}_{\text{support proper over }A}(X, \mathcal{I})$.
We conclude that $\Ker(\alpha)$ and $\Coker(\alpha)$
are in $\textit{Coh}_{\text{support proper over }A}(X, \mathcal{I})$
by Lemma \ref{lemma-systems-with-proper-support} part (1).
By induction hypothesis, more precisely because $\mathcal{K}^t$ is in $\Xi$,
we see that $\Ker(\alpha)$ and $\Coker(\alpha)$ are in
the subcategory defined in
Lemma \ref{lemma-systems-with-proper-support} part (2).
Since this is a Serre subcategory by the lemma, we conclude that the
same is true for $(\mathcal{F}_n)$ which is what we wanted to show.
\end{proof}

\bibliographystyle{plain}
\bibliography{my}
\end{document}
